\chapter{Conclusão}

\section{Sintese dos resultados}

Este trabalho teve como objetivo avaliar o desempenho de diferentes modelos na previsão da volatilidade da taxa de câmbio 
USD/BRL, comparando uma abordagem de aprendizado de máquina, representada pelo Random Forest, com modelos tradicionais da 
família GARCH, especificamente GARCH e EGARCH.

Os resultados da análise empírica evidenciaram a presença de características típicas de séries financeiras nos retornos 
cambiais, como não normalidade e dependência temporal nos retornos ao quadrado, reforçando a relevância da modelagem da 
volatilidade. Na comparação dos modelos, considerando a janela expansiva, o Random Forest apresentou os menores valores 
de MSE e QLIKE, enquanto o EGARCH obteve o menor MAE. Apesar dessas diferenças numéricas, os testes Diebold-Mariano não 
identificaram diferenças estatisticamente significativas entre os modelos para nenhuma das duas funções de perda consideradas.
Portanto, não foram encontradas evidências de superioridade preditiva estatística de um modelo em relação aos demais nessa 
especificação.

Na análise de robustez, a substituição da janela expansiva por uma janela rolante de cinco anos alterou o desempenho relativo 
dos modelos. Nessa configuração, o GARCH apresentou os melhores resultados em todas as métricas avaliadas, com diferenças 
estatisticamente significativas em relação ao Random Forest, conforme os testes Diebold-Mariano para MSE e QLIKE. Esse resultado 
demonstra que o desempenho preditivo dos modelos depende também da estratégia de estimação adotada, não sendo possível 
estabelecer uma superioridade independente dessa escolha.

\section{Contribuições}

A principal contribuição deste estudo consiste na comparação empírica entre modelos econométricos tradicionais e uma abordagem 
de aprendizado de máquina na previsão da volatilidade cambial brasileira. Além da avaliação por diferentes métricas de desempenho, 
foram utilizados testes de comparação preditiva com a correção de Harvey-Leybourne-Newbold (HLN), buscando conferir maior rigor à 
análise das diferenças entre os modelos.

Do ponto de vista aplicado, o estudo contribui para a discussão sobre a utilização de métodos de aprendizado de máquina em problemas 
de previsão financeira, especialmente em comparação com modelos econométricos amplamente utilizados na literatura. Os resultados também 
oferecem subsídios para a escolha de abordagens de previsão em aplicações relacionadas à análise de risco e ao acompanhamento da volatilidade cambial.

\section{Agenda futura}

As limitações identificadas ao longo deste estudo oferecem oportunidades para o desenvolvimento de pesquisas futuras. Uma primeira 
possibilidade consiste na utilização de medidas alternativas de volatilidade, especialmente a volatilidade realizada calculada a 
partir de dados intradiários, permitindo uma aproximação mais precisa da volatilidade latente e uma avaliação da sensibilidade dos 
resultados à proxy adotada neste trabalho.

Outra extensão seria explorar diferentes configurações do Random Forest, incluindo a seleção de variáveis, a construção de novas 
features e a frequência de reestimação do modelo. A investigação de outras abordagens de aprendizado de máquina, como o XGBoost, 
que incorporem mecanismos de dependência temporal também poderia contribuir para a comparação com os modelos econométricos tradicionais.

Além disso, seria interessante ampliar a análise para horizontes de previsão superiores a um dia, verificando se o desempenho 
relativo dos modelos se mantém em previsões de médio prazo. A incorporação de variáveis exógenas, como o índice VIX, o diferencial de 
juros entre Brasil e Estados Unidos e o CDS soberano brasileiro, também é uma possibilidade de investigação, permitindo avaliar se informações 
adicionais contribuem para melhorar a capacidade preditiva dos modelos.

Por fim, a realização de análises em diferentes subperíodos e a investigação de outras estratégias de estimação poderiam aprofundar a 
compreensão sobre a sensibilidade dos modelos às mudanças nas condições de mercado, especialmente diante de episódios de maior instabilidade 
cambial.
